\chapter{Robustness and diagnostics}
\label{ch:robust}

A kernel for a device that runs unattended has to answer two questions: what does it do
when the application breaks a rule, and how does the engineer find out what happened. The
answers in nornCore are a small number of mandatory checks, a panic record that survives
a reset, and a set of diagnostics that read the kernel's own bookkeeping.

\section{Checks and assertions}

Every public function that takes a task id, a storage pointer, a stack size or a timeout
validates it with \code{RTOS\_CHECK}, in every configuration. A violation is a
\emph{panic} with \code{RTOS\_PANIC\_ARGUMENT} and the two offending values in the record;
there is no error return for an argument that would corrupt memory, because a return
value is what the caller that made the mistake would ignore. Contracts whose violation
costs more to detect than it is worth in a release build are \code{RTOS\_ASSERT}ed
instead and compiled only with \code{RTOS\_CFG\_ASSERT}: recursive locks, the unlock
order, a lock from the idle task, an empty event mask. Chapter~\ref{ch:api} marks every
entry \emph{checked} or \emph{asserted}.

\section{Panic and the record}

\code{rtos\_panic(reason, arg1, arg2)} disables interrupts, writes a record (reason,
running task, tick, the two arguments, and for an assertion the file and line) into
memory that survives a reset, calls \code{rtos\_panicHook} and then halts
(\code{RTOS\_CFG\_PANIC\_ACTION 0}) or resets the board (1). After a reset,
\code{rtos\_init} takes the record out of that memory and \code{rtos\_getLastPanic}
returns it once; \code{rtos\_getResetCause} says what the hardware thinks happened
(power on, external, brown out, watchdog, software).

\begin{table}[htbp]
\centering\small
\begin{tabularx}{\linewidth}{@{}lX@{}}
\toprule
Reason & Meaning, \code{arg1} / \code{arg2} \\
\midrule
\code{RTOS\_PANIC\_ASSERT}         & a failed \code{RTOS\_ASSERT}; file and line in the record \\
\code{RTOS\_PANIC\_STACK\_OVERFLOW} & the canary of task \code{arg1} was overwritten (\code{STACK\_CHECK 2}) \\
\code{RTOS\_PANIC\_WATCHDOG}       & task \code{arg1} did not check in; \code{RTOS\_NO\_TASK}: the idle task \\
\code{RTOS\_PANIC\_DEADLOCK}       & task \code{arg1} would wait for a mutex held, through a chain, by itself; \code{arg2} the owner \\
\code{RTOS\_PANIC\_FAULT}          & a CPU exception; cause and address from the port \\
\code{RTOS\_PANIC\_CONFIG}         & the application's and the kernel's configuration signatures differ \\
\code{RTOS\_PANIC\_ARGUMENT}       & a mandatory check refused an argument; the values \\
\code{RTOS\_PANIC\_APPLICATION}    & \code{rtos\_panic} called by the application; its own arguments \\
\bottomrule
\end{tabularx}
\caption{Panic reasons.}
\label{tab:panic}
\end{table}

tc05 provokes four of these in turn on the boards (an application panic, a stack
overflow, a late task, a starved idle task), lets each reset the board, and checks the
record after every reset.

\section{Stack checking}

Section~\ref{sec:stacks} has the sizes; this is the machinery. With
\code{RTOS\_CFG\_STACK\_CHECK} the stack is filled with a pattern at creation and
\code{rtos\_getStackFree} counts the untouched bytes from the bottom: the high water
mark since the start. With \code{RTOS\_CFG\_STACK\_WARN} the idle task reads one task's
mark per loop and calls \code{rtos\_stackWarnHook} once for a task whose margin drops
below \code{RTOS\_PORT\_STACK\_WARN}. With \code{STACK\_CHECK 2} a canary at the bottom
is checked before every context switch and at every tick, and an overwritten canary is
\code{RTOS\_PANIC\_STACK\_OVERFLOW}. The canary catches an overflow after the fact; the
monitor warns before it.

\section{The watchdog}

\code{RTOS\_CFG\_WATCHDOG} arms the hardware watchdog (\code{rtos\_wdtInit}) and lets the
idle task kick it, but only after checking that every monitored task
(\code{rtos\_wdtMonitor}) has called \code{rtos\_wdtCheckIn} within its interval. A late
task is \code{RTOS\_PANIC\_WATCHDOG} with a record; a starved idle task, which means a
runaway task or an interrupt storm, lets the hardware reset the board without one, and
the next boot sees \code{RTOS\_RESET\_WATCHDOG}. The two cases together cover both ways
a system can hang.

\section{Deadlock detection}

A mutex wait that would close a cycle (A waits for a mutex B holds while B waits for one
A holds, through any number of tasks) is detected at the moment of the wait and is
\code{RTOS\_PANIC\_DEADLOCK}. With correct ceilings a wait never happens, so the detector
is a safety net for a wrong ceiling or an owner that blocked while holding.

\section{Run time statistics and task information}

\code{RTOS\_CFG\_RUNTIME\_STATS} measures the time each task runs, at every switch, in
counts of \code{RTOS\_RUNTIME\_HZ}. \code{rtos\_getTaskRunTime} returns a task's total,
\code{rtos\_getCpuLoad} the share of time not spent in the idle task since the previous
call. \code{rtos\_getTaskInfo} is the snapshot for a console or a debugger: state,
priority and effective priority, what the task waits on and until when, stack size and
free, run time, overruns. \code{rtos\_getOverrunCount} reads and optionally resets a
periodic task's overruns. The kernel objects also carry their own statistics: lost gives
on a semaphore, dropped sends and the high water mark on a queue, overwrites on a mailbox,
failed allocations and the low water mark on a pool, the heap's statistics.

\section{Trace}

\code{RTOS\_CFG\_TRACE} calls \code{rtos\_traceHook(event, a, b)} at every context switch
(the tasks leaving and entering), every block (task and wait kind), every wake-up, every
timer expiry and every panic, with interrupts disabled. \code{hal/trace.h} is a ring
buffer recorder to override the hook with, and the host unit tests are built on exactly
this trace: each test compares the recorded sequence of (task, event, value) with the
expected one.
